\documentclass{article}

% =========================
% Encoding & Math
% =========================
\usepackage[utf8]{inputenc}
\usepackage{amsmath, amssymb, amsfonts, amsthm}

% =========================
% Page / Graphics / Tables
% =========================
\usepackage{geometry}
\geometry{a4paper, margin=1in}

\usepackage{graphicx}
\usepackage{booktabs}
\usepackage{multirow}
\usepackage{xcolor}
\usepackage{float} % Better figure placement with [H]

% =========================
% Algorithms
% =========================
\usepackage{algorithm}
\usepackage{algpseudocode}

% Caption should be loaded after float/graphics and before hyperref in many setups
\usepackage{caption}
\usepackage{mathtools} % provides \coloneqq
% Provide a "List of Algorithms" file target (loa) cleanly
\makeatletter
\newcommand{\listofalgorithms}{%
  \section*{List of Algorithms}%
  \addcontentsline{toc}{section}{List of Algorithms}%
  \@starttoc{loa}%
}
\makeatother

% =========================
% Hyperlinks (load late)
% =========================
\usepackage{hyperref}

% =========================
% Title / Theorems / Macros
% =========================
\title{\textbf{Gradient Shake-to-Align Dual-Rank LoRA: \\ A Comprehensive Framework for Multi-Resolution Gradient Consensus}}
\author{Methodology \& Implementation Guide}
\date{\today}

\newtheorem{lemma}{Lemma}
\newtheorem{theorem}{Theorem}
\newtheorem{corollary}{Corollary}
\newtheorem{definition}{Definition}

% Utility macros
\newcommand{\norm}[1]{\left\lVert#1\right\rVert}
\newcommand{\inner}[2]{\langle #1, #2 \rangle}

% =========================
% Breakable (multi-page) algorithm environment
% =========================
\makeatletter
\newenvironment{breakablealgorithm}
  {%
   \begin{center}
   \refstepcounter{algorithm}%
   \hrule height .8pt depth 0pt \kern 2pt%
   \renewcommand{\caption}[2][\relax]{%
     {\raggedright\textbf{Algorithm \thealgorithm} ##2\par}%
     \ifx\relax##1\relax
       \addcontentsline{loa}{algorithm}{\protect\numberline{\thealgorithm}##2}%
     \else
       \addcontentsline{loa}{algorithm}{\protect\numberline{\thealgorithm}##1}%
     \fi
     \kern 2pt\hrule\kern 2pt
   }%
  }{%
   \kern 2pt\hrule\relax%
   \end{center}
  }
\makeatother

\begin{document}

\maketitle


\begin{abstract}
Parameter-efficient fine-tuning with LoRA depends critically on rank choice, yet a larger rank does not necessarily outperform a smaller one on a given dataset. This is especially true in information-limited or heterogeneous regimes, where gradients are noisy, weakly supported, or conflicting across samples and tasks. In such settings, rank selection becomes both a practical burden and an optimization problem: the key question is not simply which rank is best offline, but which rank provides the more trustworthy gradient direction at each training step under a fixed parameter budget.

We propose Gradient Shake-to-Align, a diagnostic-first framework that treats a nested rank-$r$ core and rank-$R$ extension as two gradient views of the same update. At each optimization step, cosine-based diagnostics computed from micro-batch votes measure internal consistency within each rank and directional deviation across ranks. These signals determine whether sparse cross-rank intervention is needed and which direction should be trusted. Under the nested parameterization, the low-rank head is preserved as the structural anchor, while the higher-rank branch is allowed to realize only residual semantics supported by consistency evidence. Operating in normalized gradient space, the method provides a scale-invariant way to use the internal orderliness of sample-induced gradient structure itself as a source of constraint and reference. This yields stable performance in standard single-task training and stronger robustness than baseline methods in multi-task settings with pronounced gradient conflict.


\end{abstract}


\tableofcontents
\listoffigures
\listofalgorithms

\newpage
% =========================
% NEW: Introduction + Related Work
% =========================
\section{Introduction}

\subsection{Problem Setting and Motivation}

Parameter-efficient fine-tuning (PEFT) methods such as LoRA~\cite{hu2022lora}
are built on the premise that most task-specific updates can be represented
within a low-rank subspace.
This assumption is often effective, but in practice the choice of rank is
highly data-dependent.
A larger rank does not necessarily yield better adaptation than a smaller one
on a given dataset, especially in \emph{information-limited regimes}---including
small datasets, high-entropy labels, heterogeneous supervision, or multi-task
mixtures---where gradients are noisy, weakly supported, and often inconsistent
across samples and minibatches~\cite{yu2020pcgrad,liu2021cagrad}.

This creates both a practical and a conceptual challenge.
Practically, selecting a suitable LoRA rank usually requires manual search over
discrete candidates, which is expensive and brittle.
Conceptually, existing PEFT methods mostly treat rank as a static capacity
choice, even though the reliability of the induced gradients can vary over the
course of training.
A higher-rank parameterization provides more flexibility, but may absorb
stochastic noise or task-specific interference.
A lower-rank parameterization is more restrictive, yet can sometimes expose a
cleaner and more stable optimization direction.

These observations suggest that the central question is not simply
\emph{which rank is better}, but rather:
\emph{at each optimization step, which rank provides the more trustworthy
gradient direction for the current data and training state?}
If such reliability can be inferred online, rank selection no longer needs to
be treated purely as an offline hyperparameter search problem.

\subsection{Core Insight: Nested Ranks as Online Reliability Probes}

Our key insight is that nested LoRA parameterizations with different ranks
provide complementary resolutions of the same update signal.
The smaller rank captures the most compressible and structurally stable part of
the adaptation, while the larger rank retains additional flexibility for
residual semantics that cannot be expressed in the low-rank core.

Importantly, we do not assume that the larger rank should always dominate.
Instead, the interaction between the two ranks provides an intrinsic diagnostic
signal.
When low- and high-rank views exhibit consistent gradient structure, the update
is more likely to be reliable and semantically aligned across resolutions.
When they disagree, the discrepancy may reveal high-rank overfitting,
low-rank under-capacity, or gradient conflict induced by noisy or heterogeneous
samples.

This turns nested rank design into more than a capacity decomposition:
it becomes a mechanism for online reliability estimation.
Rather than manually searching for a single best rank, we use the internal
structure of gradient space to determine which direction is more trustworthy at
each step, and to constrain high-rank execution accordingly.

\subsection{Gradient Shake-to-Align: Reliability-Guided Direction Selection}

We introduce \emph{Gradient Shake-to-Align}, a diagnostic-first framework
that uses nested low- and high-rank LoRA views to identify trustworthy update
directions during training.
At each optimization step, gradients are partitioned into multiple voters,
allowing us to measure internal consistency within each rank and directional
agreement across ranks.

These signals are used not to enforce constant alignment, but to determine
whether intervention is necessary.
When the gradient structure is coherent, each rank proceeds with its own update.
When inconsistency or conflict becomes prominent, the method activates a sparse
cross-rank reconciliation mechanism that selects the more trustworthy direction
based on consistency evidence.
The low-rank head is preserved as the structural anchor, while the higher-rank
component is constrained to realize only residual semantics outside the trusted
core.

In this sense, our method uses the geometry and orderliness of the
sample-induced gradient space itself as a source of constraint and reference.
It improves optimization stability in standard single-task settings and is
particularly beneficial in multi-task regimes, where gradient conflict is more
severe and direction reliability becomes a central challenge.

\subsection{Contributions}

We summarize our contributions as follows:

\begin{itemize}

    \item \textbf{We reformulate rank selection as an online reliability problem.}
    Instead of assuming that a larger LoRA rank is always preferable, or
    treating rank choice purely as offline hyperparameter tuning, we ask which
    rank provides the more trustworthy gradient direction at each training step.

    \item \textbf{We propose a nested-rank diagnostic framework for trustworthy update selection.}
    Our method uses low- and high-rank LoRA parameterizations as complementary
    gradient views and leverages their internal consistency and cross-rank
    agreement to evaluate directional reliability online.

    \item \textbf{We introduce a structure-driven execution rule under a fixed parameter budget.}
    Without exceeding the trainable parameter count of a comparable
    single large-rank LoRA, our framework preserves the low-rank head as a
    trusted structural core and constrains the high-rank component to realize
    residual semantics only when supported by consistency evidence.

    \item \textbf{We provide a gradient-space regularization perspective that is especially effective under conflict.}
    The method uses the orderliness of sample-induced gradient structure itself
    as a source of constraint and reference, yielding stable performance in
    standard single-task training and stronger robustness than baseline methods
    in multi-task settings with pronounced gradient conflict.
\end{itemize}
\section{Related Work and Limitations}

\subsection{Gradient Conflict, Agreement, and Diagnostic Signals}

A line of work in multi-task and multi-objective optimization studies how to handle
gradient conflicts when different objectives or data sources disagree.
Representative approaches include PCGrad~\cite{yu2020pcgrad},
which projects conflicting gradients to reduce interference,
and CAGrad~\cite{liu2021cagrad}, which seeks a conflict-averse descent direction
by optimizing a worst-case objective so that progress on one task does not come
at the expense of severe degradation on others.
These methods highlight the limitations of naive gradient averaging
in the presence of heterogeneous or conflicting signals.

However, these methods primarily aim to construct a better descent direction
once conflict is observed, rather than asking why the conflict arises or whether
one parameterization is more trustworthy than another.
In practice, disagreement may stem from stochastic noise~\cite{yin2018gradient},
capacity mismatch, or structured but weakly supported signals.
Treating all conflicts uniformly can therefore obscure whether disagreement
reflects noise, insufficient rank, or meaningful heterogeneity.

Our work departs from immediate conflict resolution and instead treats
cross-rank gradient agreement as an \emph{online reliability signal}.
Rather than modifying the descent objective directly,
we compare internal consistencies across nested parameter resolutions
and use this contrast to determine \emph{whether}, \emph{when}, and
\emph{which rank direction} should guide cross-resolution intervention.

\subsection{Low-Rank Adaptation, Adaptive Capacity, and Subspace Drift}

LoRA~\cite{hu2022lora} parameterizes task-specific updates in a low-rank subspace
and has become a standard baseline for parameter-efficient fine-tuning.
Subsequent methods, including adaptive and dynamic rank variants such as
AdaLoRA~\cite{zhang2023adalora}, allocate capacity based on parameter magnitudes
or heuristic importance measures.
More recent work explores alternative parameterizations to improve stability
and expressiveness, including DoRA~\cite{liu2024dora} and spectral approaches~\cite{wang2024svft}.

Despite these advances, most PEFT methods implicitly assume that the effective
adaptation subspace is sufficiently stable throughout training,
or that capacity reallocation alone is sufficient to address mismatch.
In contrast, empirical evidence suggests that the effective update subspace
can evolve over time~\cite{aghajanyan2021intrinsic},
particularly in low-data or high-entropy regimes.
This mismatch can lead to unstable updates when low-rank structure
is trusted prematurely or enforced without reliability evidence.

Our framework addresses this issue from a \emph{diagnostic perspective}.
Rather than adapting capacity solely based on parameter statistics,
we use cross-rank gradient consistency as an intrinsic reliability estimator.
Low-rank structure is preserved when supported by internal agreement,
while higher-rank flexibility remains available but constrained when signals are
unstable.
This shifts the focus from static rank control to evidence-driven execution
control under a fixed parameter budget.

\subsection{Residual Information and Evidence-Driven Consolidation}

Work on gradient compression and approximation shows that discarding update
components without retaining residual information can impair convergence, and
that error feedback is often important for stability~\cite{karimireddy2019error}.
Low-rank approximation methods such as PowerSGD~\cite{vogels2019powersgd}
further illustrate that structured update compression can be effective.

Our setting is different: we are not primarily compressing gradients for
communication or storage.
The relevance of this line of work is conceptual.
It suggests that information outside a compressed core should not be discarded
blindly.
Our method follows the same principle, but uses cross-rank agreement to decide
when residual semantics should remain decoupled and when alignment is justified.

\subsection{Summary of Limitations in Existing Paradigms}

Across these lines of work, existing paradigms tend to
(i) resolve gradient conflicts without diagnosing their origin,
(ii) impose or adapt low-rank structure without an explicit reliability criterion
grounded in gradient agreement,
or (iii) treat residual information uniformly without separating noise
from structured but weakly supported signals.

Our approach bridges these gaps by leveraging cross-rank gradient consistency
as an intrinsic online reliability signal.
This enables sparse and targeted intervention in nested-rank training,
reframing PEFT as an evidence-driven execution process rather than a static
structural constraint or a one-shot rank-selection problem.
% =========================
% Existing section (keep your original text)
% =========================
% Method Overview (Component-first, Algorithm-ordered)
% =========================
\section{Method}
\label{sec:method}


\subsection{Overview}
\label{sec:method-overview}

We propose \textit{Gradient Shake-to-Align}, a diagnostic-first framework for Dual-Rank LoRA.
Our design follows a single guiding principle:

\begin{quote}
\textbf{Diagnosis before intervention.}
We do not enforce low-rank behavior \emph{a priori}; instead, we use online evidence from
multi-view gradient consistency to decide \emph{when} cross-resolution alignment
is necessary during training.
\end{quote}

\subsubsection{Objects and Parameterization (What exists in the system)}
For each LoRA block $b$, we maintain a nested Dual-Rank update of target rank $R$
with a stable core rank $r<R$:
\begin{equation}
\Delta W_{R,b}
=
\underbrace{U_{r,b}V_{r,b}}_{\text{core (rank } r\text{)}}
+
\underbrace{U_{hi,b}V_{hi,b}}_{\text{extension (rank } R-r\text{)}}.
\end{equation}

\paragraph{Initialization (part of the method, not a separate chapter).}
We initialize the extension to contribute zero at $t=0$:
\begin{equation}
U_{r,b},V_{r,b}\sim \text{Kaiming},\qquad
U_{hi,b}\sim \mathcal{N}(0,\sigma^2),\qquad
V_{hi,b}=0,
\end{equation}
so that $\Delta W_{R,b}(0)=U_{r,b}V_{r,b}$.
As a result, the high-resolution branch starts as a pure \emph{reservoir}
that can absorb heterogeneous or unreliable gradients
without immediately affecting the low-rank core.

\subsubsection{Diagnostic Evidence (How we look at the world)}
\begin{itemize}
    \item \textbf{Votes per step (implementation-aligned):}
    at each training step we collect gradient votes by partitioning each task minibatch into
    microbatch windows of size $s_{\mathrm{pv}}$ (``samples per vote'').
    The effective number of votes at step $t$ is therefore
    \[
      V_t \coloneqq \sum_{m\in\mathcal{M}} \left\lceil \frac{|B_m|}{s_{\mathrm{pv}}} \right\rceil,
    \]
    which is constant only when task batch sizes are fixed.
    \item \textbf{Two diagnostic views:}
    \begin{itemize}
        \item \textbf{$r$-view:} gradients restricted to $(U_r,V_r)$,
        zero-padded to rank-$R$ space.
        \item \textbf{$R$-view:} gradients over $(U_r,V_r,U_{hi},V_{hi})$.
    \end{itemize}
\end{itemize}

From these votes, we compute per-step diagnostic statistics:
\begin{itemize}
    \item \textbf{Internal consistencies:} $C_{r,b}$ and $C_{R,b}$,
    measuring agreement among votes within each view.
    \item \textbf{Cross-resolution retention ratio:} $A_b$,
    measuring how much of the aggregated rank-$R$ update remains explainable by the rank-$r$ subspace.
\end{itemize}
All statistics are computed directly from the microbatch votes within the current optimization step.
Precise definitions are introduced where needed.

\subsubsection{Direction Construction (Per-view, mean-based)}
In the experimental setting considered in this paper,
each view constructs its execution direction using the
mean of its corresponding gradient votes.
This choice provides a simple and stable reference direction
while allowing the diagnostic statistics to determine
whether further cross-resolution alignment is required.

\subsubsection{Global Trigger (Gate-0): When to reconcile across resolutions}
Cross-resolution reconciliation is designed to be \emph{sparse}.
A global trigger (Gate-0) activates reconciliation only when
there is evidence of unreliability or disagreement,
as indicated by low internal consistency in either view
or large directional deviation between the two views.
When Gate-0 does not trigger, updates are applied independently
within each view.

\subsubsection{Cross-Resolution Reconciliation (Trusted reference + constrained execution)}
When Gate-0 triggers, we identify a trusted reference view
based on relative internal consistency.
Execution is then performed under the nested-parameterization constraint:
\textbf{the rank-$R$ update is not allowed to overwrite the rank-$r$ core}.
Accordingly, the low-rank core is updated using the $r$-view direction,
while the high-rank extension realizes any residual semantics
required to maintain consistency with the $R$-view.

This constrained execution rule ensures semantic closure
under the nested parameterization,
allowing higher-resolution information to assist training
without destabilizing the low-rank representation.

\subsection{Part A: Gradient Shake-to-Align (Per-step Diagnosis and Execution)}
\label{sec:partA}
\subsubsection{Core Methodology: Ordered Pipeline}
\label{sec:core-methodology}
We now describe the end-to-end per-step pipeline in execution order.

\paragraph{Step 1: Build Diagnostic Statistics from Votes}
At each training step, we construct a set of gradient votes for each LoRA block
via a memory-safe \emph{serial microbatch voting} procedure:
we iterate over microbatch windows, run \texttt{forward}+\texttt{backward} per window,
and store the resulting (scale-normalized) LoRA gradients as votes.
We then compute per-step diagnostic statistics purely from the current microbatch votes:
\[
C_{r,b},\quad C_{R,b},\quad A_b,
\]
where $C_{r,b}$ and $C_{R,b}$ measure instantaneous internal agreement among votes
within the $r$- and $R$-views respectively, and $A_b$ measures
cross-resolution retention of the aggregated rank-$R$ update by the rank-$r$ subspace.

\paragraph{Step 2: Mean Direction Construction (Per-view)}
For each view, we construct a representative execution direction
using the mean of its corresponding gradient votes:
\[
\bar v_r := \frac{1}{V_t}\sum_{k=1}^{V_t} v^{(k)}_{r,b},
\qquad
\bar v_R := \frac{1}{V_t}\sum_{k=1}^{V_t} v^{(k)}_{R,b}.
\]
These mean directions serve as stable references,
while all diagnostic decisions are made using the underlying vote statistics.
All mean directions and subsequent pulling operations are interpreted in the scale-invariant (normalized) gradient space.

\paragraph{Step 3: Gate-0 (Global Trigger): When to reconcile across resolutions}
Cross-resolution reconciliation is designed to be \emph{sparse}.
A global trigger (Gate-0) activates reconciliation only when
there is evidence of unreliability or disagreement.
Specifically, we enter reconciliation if either:
(i) at least one view exhibits high internal unreliability under a consistency-derived noise score, or
(ii) the two views disagree strongly under mean directions.

\paragraph{(i) Noise-summary trigger from internal consistency.}
\begin{equation}
\label{eq:noise-gates}
b_{r,b} := \sigma\!\left(\kappa_N(\tau_{\mathrm{noise}} - C_{r,b})\right),
\qquad
b_{R,b} := \sigma\!\left(\kappa_N(\tau_{\mathrm{noise}} - C_{R,b})\right),
\end{equation}
\begin{equation}
\label{eq:noise-summary}
N_b := \max(b_{r,b},\, b_{R,b}).
\end{equation}

\paragraph{(ii) Directional deviation trigger.}
\begin{equation}
\label{eq:direction-deviation-raw}
D_b := 1 - \mathrm{CosSim}\!\left(\bar v_r,\; \mathrm{head}(\bar v_R)\right).
\end{equation}
We use $\mathrm{head}(\bar v_R)$ to denote the first $r$ coordinates
of the rank-$R$ direction under the nested parameterization.
This is not identical to the $r$-view direction $\bar v_r$.

\paragraph{Gate-0 (OR logic).}
\begin{equation}
\label{eq:gate0-final}
b \in \mathcal{T}
\;\;\Longleftrightarrow\;\;
\Big(N_b \ge \tau_N\Big)
\;\;\lor\;\;
\Big(D_b \ge \tau_D\Big).
\end{equation}

\noindent
If $b\notin\mathcal{T}$, updates are applied independently in each view
using mean directions.
If $b\in\mathcal{T}$, we perform cross-resolution reconciliation
as described below.

\paragraph{Remark: Diagnostic statistics vs.\ execution directions.}
We emphasize that diagnostic statistics $(C_{r,b}, C_{R,b}, A_b)$
are computed solely from the vote distributions and are independent of
the choice of execution direction.
Here, $A_b$ is a retention-style diagnostic rather than a primary directional gating variable.
Mean directions act only as representatives for execution and do not
alter the underlying diagnostic evidence.

\paragraph{Step 4: Trusted Reference Selection and Pull Strength}
When reconciliation is triggered, we quantify the relative reliability
contrast between the two views by
\[
\Delta C_b := C_{r,b} - C_{R,b}.
\]

Rather than performing a hard switch, we interpret $\Delta C_b$ as defining:
(i) a \emph{trusted reference view}, and
(ii) a continuous \emph{pull strength} controlling how strongly the less
reliable view is aligned to the reference.

\paragraph{Reference direction.}
We select the trusted head direction as
\[
g_{\mathrm{ref}} :=
\begin{cases}
\bar v_r, & \Delta C_b \ge 0 \quad (\text{$r$ view more reliable}), \\
\mathrm{head}(\bar v_R), & \Delta C_b < 0 \quad (\text{$R$ view more reliable}).
\end{cases}
\]

\paragraph{Pull strength.}
The magnitude of the reliability gap is converted into a pull coefficient:
\[
\lambda_b := \mathrm{clip}\!\left(\frac{|\Delta C_b|-\delta}{1-\delta},\,0,\,1\right),
\]
so that small gaps induce negligible alignment,
while large gaps result in stronger reconciliation.

\paragraph{Soft directional pulling.}
Let $\bar v_R=(\bar v_{R,h},\bar v_{R,t})$ denote the head--tail decomposition
of the rank-$R$ mean direction.
The pulled head targets are defined as
\[
(g_h^\dagger,\; g_{R,h}^\dagger) :=
\begin{cases}
\big(\bar v_r,\; (1-\lambda_b)\bar v_{R,h} + \lambda_b \bar v_r\big), & \Delta C_b \ge 0, \\[6pt]
\big((1-\lambda_b)\bar v_r + \lambda_b \bar v_{R,h},\; \bar v_{R,h}\big), & \Delta C_b < 0.
\end{cases}
\]
Only the less reliable view is interpolated toward the reference,
while the trusted view remains unchanged.

\paragraph{Step 5: Constrained Execution (Semantic Closure under Nesting)}

Execution is performed in a scale-normalized gradient space
to ensure invariance to LoRA scaling factors. All directions are now interpreted in normalized space.
Each branch applies scaling
\[
\Delta W_r = s_r B_r A_r,
\qquad
\Delta W_h = s_h B_h A_h,
\qquad
\Delta W_R = \Delta W_r + \Delta W_h.
\]

Before any execution or compensation, we define the scale-invariant (normalized) gradients:
\[
g_{A_r} := g_{A_r}/s_r,\quad
g_{B_r} := g_{B_r}/s_r,
\qquad
g_{A_h} := g_{A_h}/s_h,\quad
g_{B_h} := g_{B_h}/s_h.
\]
All reconciliation, pulling, and compensation are performed in this scale-invariant space.


\paragraph{Why scale-invariant reconciliation?}
LoRA scaling factors (e.g., shared $\alpha$ conventions)
introduce magnitude differences that do not reflect
intrinsic gradient reliability.
Performing diagnostics and reconciliation in a normalized space
ensures that cosine-based consistency statistics
and cross-resolution pulling depend only on directional structure,
not on arbitrary scaling choices.


\paragraph{Head execution.}
All pulling targets are defined in the scale-invariant space.
Under the nested dual-rank parameterization, the rank-$R$ update is not allowed to overwrite the rank-$r$ core.
We therefore execute the head strictly using the trusted head direction:

\[
g_r^{\mathrm{exec}} \coloneqq g_h^\dagger.
\]
where $g_r^{exec}$ denotes the concatenated executed gradient on $(A_r,B_r)$.

\paragraph{Tail realization via residual compensation.}
Let the executed rank-$R$ direction be
\[
g_R^{exec} = (g_{R,h}^\dagger,\; \bar v_{R,t}),
\]
where $\bar v_{R,t}$ is the tail (extension) component of the mean direction.

Define the head discrepancy:
\[
r_{\mathrm{head}} =
\mathrm{head}(g_R^{exec}) - g_r^{exec}.
\]

Mapping this discrepancy back to weight space yields a first-order residual:
\[
\Delta W_{\mathrm{res}}
\approx
B_r\,\delta A_r + \delta B_r\,A_r.
\]

To preserve rank-$r$ execution while realizing the full rank-$R$ semantics,
we solve for extension updates $(\delta A_h,\delta B_h)$ such that
\[
B_h\,\delta A_h + \delta B_h\,A_h
\approx
\Delta W_{\mathrm{res}}.
\]

We use ridge-regularized least squares:
\[
\delta B_h^*
=
\arg\min_X
\|X A_h - \Delta W_{\mathrm{res}}\|_F^2
+
\lambda_{\mathrm{ls}}\|X\|_F^2,
\qquad
\delta A_h^*
=
\arg\min_Y
\|B_h Y - \Delta W_{\mathrm{res}}\|_F^2
+
\lambda_{\mathrm{ls}}\|Y\|_F^2,
\]
with closed-form solutions
\[
\delta B_h^*
=
\Delta W_{\mathrm{res}}
A_h^\top
(A_h A_h^\top + \lambda_{\mathrm{ls}} I)^{-1},
\qquad
\delta A_h^*
=
(B_h^\top B_h + \lambda_{\mathrm{ls}} I)^{-1}
B_h^\top
\Delta W_{\mathrm{res}}.
\]

The executed extension gradient becomes
\[
g_{hi}^{exec}
=
\bar v_{R,t}
+
\mathrm{vec}(\delta A_h^*, \delta B_h^*).
\]

\paragraph{Justification of the first-order residual approximation.}
The residual mapping
\(
\Delta W_{\mathrm{res}}
\approx
B_r\,\delta A_r + \delta B_r\,A_r
\)
follows from the first-order expansion of the bilinear LoRA parameterization.
Under standard small-step assumptions of first-order optimizers (e.g., AdamW),
the update of $(A_r,B_r)$ remains in a locally linear regime,
making the ridge-based least-squares compensation a consistent approximation
of the residual weight mismatch.

\paragraph{Write-back to optimizer.}
Finally, gradients are restored to the original scaled parameter space:
\[
g_{A_r}^{write}=s_r\,g_{A_r}^{exec},\quad
g_{B_r}^{write}=s_r\,g_{B_r}^{exec},
\qquad
g_{A_h}^{write}=s_h\,g_{A_h}^{exec},\quad
g_{B_h}^{write}=s_h\,g_{B_h}^{exec}.
\]

This construction guarantees:
(i) exact preservation of the rank-$r$ core,
(ii) semantic closure of the rank-$R$ update,
(iii) invariance to LoRA scaling choices.
Moreover, as the head discrepancy $r_{\mathrm{head}}$ shrinks (e.g., when $\lambda_b$ increases),
the residual $\Delta W_{\mathrm{res}}$ approaches zero, and the LS compensation
naturally vanishes.


\subsubsection{Limiting Regimes (Explanatory)}

The proposed mechanism is continuous and evidence-driven.
The following familiar cases arise only as limiting patterns
of the reliability contrast $\Delta C_b$ and deviation signal $D_b$:

\begin{itemize}

    \item \textbf{Low reliability regime.}
    When one or both views exhibit high internal unreliability
    ($N_b \ge \tau_N$),
    reconciliation is triggered to prevent unstable execution.
    In this regime, intervention is driven by unreliability
    rather than directional disagreement.

    \item \textbf{Capacity insufficiency regime.}
    When $C_{R,b} \gg C_{r,b}$,
    the high-resolution view demonstrates stronger internal agreement.
    The framework interprets this as evidence that the rank-$r$
    core lacks sufficient capacity,
    and softly aligns the $r$-view toward the trusted $R$-head direction.

    \item \textbf{High-frequency overfitting regime.}
    When $C_{r,b} \gg C_{R,b}$,
    the low-rank core is internally stable while the high-rank extension
    ex exhibits lower consistency.
    The mechanism therefore regularizes the $R$-view head toward
    the $r$-view direction, suppressing unreliable high-frequency components.

    \item \textbf{Stable directional disagreement.}
    When both views are internally consistent
    yet directionally misaligned ($D_b$ large),
    the framework performs soft pulling without hard switching.
    This corresponds to structured but resolution-dependent semantics,
    reconciled through nested execution rather than constraint override.

\end{itemize}



\subsubsection{Computational Complexity and Implementation Notes}
\label{sec:complexity}

\paragraph{Algorithmic complexity.}
Under a fixed global batch size,
the proposed framework does not change the order of forward or backward
computation with respect to baseline LoRA.
All gradient votes are computed within the same optimization step,
and reconciliation operates on low-rank parameter blocks.
Therefore, the overall computational complexity
remains of the same asymptotic order as standard PEFT training.

\paragraph{Vote-level gradient extraction.}
The primary overhead arises from obtaining vote-level gradients
for consistency diagnostics.
In our current implementation, we adopt a memory-safe
serial microbatch strategy within each step.
This introduces a constant-factor slowdown relative to baseline.

\paragraph{Parallel implementation potential.}
Algorithmically, vote extraction can be parallelized
using per-sample or per-vote gradient computation (e.g., vectorized Jacobian methods),
which allows the diagnostic statistics to be computed
without repeated backward passes.
In such implementations, the practical overhead can be substantially reduced,
making the runtime comparable to baseline training up to constant factors.


\section{Algorithm}

\paragraph{Multi-task iteration (code-matched).}
In the multi-task setting, each global step samples exactly one minibatch from \emph{each} task dataloader.
Thus a single optimizer step processes the set $\{B_m\}_{m=1}^{M}$ simultaneously, where each $B_m$
is independently shuffled within its own task stream.
Votes are then extracted by microbatching each $B_m$ into vote windows.

\paragraph{Multi-task weighting (implementation-aligned).}
Let $|B_m^{(j)}|$ denote the number of samples in window $j$ of task $m$ and $N_t \coloneqq \sum_{m\in\mathcal{M}} |B_m|$ be the total number of samples across tasks at step $t$.
During vote construction we weight each window loss/gradient by $w_{m,j} \coloneqq |B_m^{(j)}|/N_t$ when accumulating the step loss and the optimizer gradients.


\begin{breakablealgorithm}
\caption{Gradient Shake-to-Align (Single-Timescale, Scale-Invariant Routing)}
\begin{algorithmic}[1]

\Require Dual ranks $(r,R)$; voter set $\mathcal{V}$ with $V_t:=|\mathcal{V}|$.
\Require Noise-gate parameters $(\tau_{\mathrm{noise}}, \kappa_N, \tau_N)$; directional threshold $\tau_D$; routing margin $\delta$; ridge $\lambda_{\mathrm{ls}}$.
\Require LoRA scaling factors $(s_r,s_h)$ for core and extension branches.
\Statex \textit{Note: All $v^{(k)}$, $\bar v_\chi$, and $g_\cdot$ inside the algorithm denote normalized gradients.}

\While{training step $t=1,2,\dots$}

    \State \textbf{Serial microbatch evidence construction (implementation-aligned):}
    \State Obtain voter gradients $\{v^{(k)}_{r,b}, v^{(k)}_{R,b}\}_{k\in\mathcal{V}}$ for each LoRA block $b$.
    \Comment{$\mathcal{V}$ indexes microbatch windows; in multi-task it spans windows across all tasks.}

    \State \textbf{Scale-normalize voter gradients:}
    \State Convert parameter gradients to scale-invariant space by dividing by branch scales $(s_r,s_h)$.

    \State \textbf{Diagnostic computation:}
    \State Compute instantaneous statistics $(C_{r,b}, C_{R,b}, A_b)$ from normalized voters,
    where $A_b$ denotes a cross-resolution retention ratio.

    \For{each block $b$}

        \State \textbf{Mean directions (execution representatives):}
        \State \quad $\bar v_r \gets \frac{1}{V_t}\sum_{k\in\mathcal{V}} v^{(k)}_{r,b}$
        \State \quad $\bar v_R \gets \frac{1}{V_t}\sum_{k\in\mathcal{V}} v^{(k)}_{R,b}$

        \State \textbf{Gate-0 signals (cheap):}
        \State \quad $b_{r,b} \gets \sigma\!\left(\kappa_N(\tau_{\mathrm{noise}} - C_{r,b})\right)$
        \State \quad $b_{R,b} \gets \sigma\!\left(\kappa_N(\tau_{\mathrm{noise}} - C_{R,b})\right)$
        \State \quad $N_b \gets \max(b_{r,b}, b_{R,b})$
        \State \quad $D_b \gets 1-\mathrm{CosSim}(\bar v_r,\mathrm{head}(\bar v_R))$

        \State \textbf{Gate-0 (global trigger; cheap OR):}
        \State \quad $\mathcal{T}_b \gets [N_b \ge \tau_N] \lor [D_b \ge \tau_D]$

        \If{$\mathcal{T}_b$}

            \State \textbf{Reliability contrast and pull strength:}
            \State \quad $\Delta C_b \gets C_{r,b}-C_{R,b}$
            \State \quad $\lambda_b \gets \mathrm{clip}\!\left(\frac{|\Delta C_b|-\delta}{1-\delta},\,0,\,1\right)$

            \State \textbf{Soft cross-resolution pulling (head level):}
            \State \quad \textbf{if} $\Delta C_b \ge 0$ \textbf{then} \Comment{$r$ view more reliable}
            \State \qquad $g_h^\dagger \gets \bar v_r$
            \State \qquad $g_{R,h}^\dagger \gets (1-\lambda_b)\mathrm{head}(\bar v_R) + \lambda_b \bar v_r$
            \State \quad \textbf{else} \Comment{$R$ view more reliable}
            \State \qquad $g_h^\dagger \gets (1-\lambda_b)\bar v_r + \lambda_b \mathrm{head}(\bar v_R)$
            \State \qquad $g_{R,h}^\dagger \gets \mathrm{head}(\bar v_R)$
            \State \quad \textbf{end if}

            \State \textbf{Constrained execution under nesting (with residual compensation):}
            \State \quad $g_r \gets g_h^\dagger$
            \State \quad Compute head residual in weight space $\Delta W_{\mathrm{res}}$ from $(g_{R,h}^\dagger - g_h^\dagger)$.
            \State \quad Solve ridge LS compensation on extension branch with $\lambda_{\mathrm{ls}}$ to realize $\Delta W_{\mathrm{res}}$.
            \State \quad $g_{hi} \gets \mathrm{tail}(\bar v_R) + \mathrm{Compensate}(\Delta W_{\mathrm{res}};\lambda_{\mathrm{ls}})$

        \Else

            \State \textbf{Default execution (no reconciliation; mean-only):}
            \State \quad $(g_r,g_{hi}) \gets (\bar v_r,\mathrm{tail}(\bar v_R))$

        \EndIf

        \State \textbf{Rescale and write-back:}
        \State Multiply executed gradients by $(s_r,s_h)$ before writing to optimizer buffers.

    \EndFor

    \State \texttt{Optimizer.step()}

\EndWhile
\end{algorithmic}
\end{breakablealgorithm}

\begin{breakablealgorithm}
\caption{Multi-Task Voter Construction (Implementation-Aligned Serial Microbatching)}
\begin{algorithmic}[1]
\Require Task set $\mathcal{M}=\{1,\dots,M\}$; at global step $t$ we draw one minibatch per task,
i.e., $\{B_m\}_{m\in\mathcal{M}}$ with $B_m \sim \mathcal{D}_m$ (via independent task dataloader iterators).
\Require Vote microbatch size $s_{\mathrm{pv}}$ (samples per vote); allow-tail flag.
\Ensure Voter set $\mathcal{V}$ and voter gradients $\{v^{(k)}_{r,b}, v^{(k)}_{R,b}\}_{k\in\mathcal{V}}$ for each block $b$.

\State Initialize $\mathcal{V}\gets\emptyset$.

\For{each task $m\in\mathcal{M}$}
    \State \textbf{Window-level voters per task (implementation-aligned):}
    \State Partition $B_m$ into windows of size $s_{\mathrm{pv}}$: $\{B_m^{(j)}\}_{j=1}^{\lceil |B_m|/s_{\mathrm{pv}}\rceil}$.
    \For{$j=1,\dots,\lceil |B_m|/s_{\mathrm{pv}}\rceil$}
        \State Compute gradient votes for each block $b$:
        \State \quad $v^{(m,j)}_{r,b} \gets \nabla_{(U_r,V_r)} \mathcal{L}_m(B_m^{(j)})$
        \State \quad $v^{(m,j)}_{R,b} \gets \nabla_{(U_r,V_r,U_{hi},V_{hi})} \mathcal{L}_m(B_m^{(j)})$
        \State Add voter index $(m,j)$ to $\mathcal{V}$.
    \EndFor
\EndFor

\State \textbf{Return} $\mathcal{V}$ and voter gradients to Algorithm~1.
\Comment{Algorithm~1 uses $\mathcal{V}$ uniformly: $\bar v_\chi=\frac{1}{|\mathcal{V}|}\sum_{k\in\mathcal{V}} v^{(k)}_{\chi,b}$.}

\end{algorithmic}
\end{breakablealgorithm}
\section{Experimental Results}
\label{sec:exp-results}

\subsection{Multi-Task Results}

For the multi-task RoBERTa study, the main table reports validation
\textit{final} accuracy for the official \textit{mix3} and \textit{mix4}
final runs. We compare \textit{baseline r}, \textit{baseline R},
\textit{CAGrad r}, \textit{CAGrad R}, \textit{ours r}, and \textit{ours R}.
The method-family averages are
$\frac{1}{2}(\textit{r}+\textit{R})$ inside each family.

\begin{table}[H]
\centering
\caption{Multi-task RoBERTa validation \textit{final} accuracy.}
\label{tab:multi-final}
\resizebox{\textwidth}{!}{%
\begin{tabular}{lccccccccc}
\toprule
Mix & baseline r & baseline R & CAGrad r & CAGrad R & ours r & ours R & Baseline Avg & CAGrad Avg & Ours Avg \\
\midrule
mix3 & 0.8148 & 0.7985 & 0.8129 & 0.8116 & \textbf{0.8183} & 0.8174 & 0.8066 & 0.8122 & \textbf{0.8179} \\
mix4 & 0.8381 & 0.8388 & 0.8308 & 0.8328 & \textbf{0.8403} & 0.8401 & 0.8385 & 0.8318 & \textbf{0.8402} \\
mix3+mix4 & 0.8265 & 0.8186 & 0.8218 & 0.8222 & \textbf{0.8293} & 0.8288 & 0.8226 & 0.8220 & \textbf{0.8290} \\
\bottomrule
\end{tabular}%
}
\end{table}

\begin{table}[H]
\centering
\caption{Mean seed variance over the retained multi-task evaluation window.
At each kept evaluation point we compute variance across random seeds for the
current \textit{final}, the running \textit{maximum}, and
\textit{total}=0.5\cdot\textit{final}+0.5\cdot\textit{maximum}, then average
those variances over the window. Lower is better.}
\label{tab:multi-seed-variance}
\resizebox{\textwidth}{!}{%
\begin{tabular}{lccccccccc}
\toprule
\multirow{2}{*}{Mix} & \multicolumn{3}{c}{Baseline Avg} & \multicolumn{3}{c}{CAGrad Avg} & \multicolumn{3}{c}{Ours Avg} \\
\cmidrule(lr){2-4}\cmidrule(lr){5-7}\cmidrule(lr){8-10}
 & Final Var & Max Var & Total Var & Final Var & Max Var & Total Var & Final Var & Max Var & Total Var \\
\midrule
mix3 & \textbf{0.000152} & \textbf{0.000145} & \textbf{0.000144} & 0.000199 & 0.000197 & 0.000197 & 0.000177 & 0.000146 & 0.000154 \\
mix4 & 0.000074 & 0.000068 & 0.000070 & \textbf{0.000043} & \textbf{0.000044} & \textbf{0.000043} & 0.000272 & 0.000253 & 0.000259 \\
Overall & \textbf{0.000113} & \textbf{0.000106} & \textbf{0.000107} & 0.000121 & 0.000121 & 0.000120 & 0.000224 & 0.000200 & 0.000207 \\
\bottomrule
\end{tabular}%
}
\end{table}

Table~\ref{tab:multi-final} shows that the proposed method leads on the
\textit{final} validation accuracy for both mixtures and for their average.
Table~\ref{tab:multi-seed-variance} shows the corresponding window-mean
seed-variance view.

\begin{table}[H]
\centering
\caption{Worst-case \textit{final} accuracy in the multi-task study, defined as
$\min(r, R)$ inside each method family.}
\label{tab:multi-worst-final}
\begin{tabular}{lcccc}
\toprule
Mix & Baseline Worst & CAGrad Worst & Ours Worst & Margin over Baseline & Margin over CAGrad \\
\midrule
mix3 & 0.7985 & 0.8116 & \textbf{0.8174} & \textbf{0.0189} & \textbf{0.0058} \\
mix4 & 0.8381 & 0.8308 & \textbf{0.8401} & \textbf{0.0020} & \textbf{0.0093} \\
Overall & 0.8183 & 0.8212 & \textbf{0.8288} & \textbf{0.0105} & \textbf{0.0076} \\
\bottomrule
\end{tabular}
\end{table}

\subsection{Single-Task Results}

For the single-task study, we fix one shared epoch budget per task across both
backbones: RTE uses epoch~7, MRPC uses epoch~7, and SST-2 uses epoch~2.  For
the main tables below we report validation \textit{final} accuracy. We compare
\textit{baseline r}, \textit{baseline R}, \textit{ours r}, and \textit{ours R}.
The task-level average compares
$\frac{1}{2}(\textit{ours r final}+\textit{ours R final})$ against
$\frac{1}{2}(\textit{baseline r final}+\textit{baseline R final})$.

\begin{table}[H]
\centering
\caption{RTE validation \textit{final} accuracy with a shared epoch-7 budget.}
\label{tab:single-rte-fixed}
\small
\begin{tabular}{lcccc}
\toprule
\multirow{2}{*}{Backbone} & baseline r & baseline R & ours r & ours R \\
\midrule
DeBERTa & 0.8383 & 0.8318 & \textbf{0.8390} & \textbf{0.8390} \\
RoBERTa & 0.7444 & \textbf{0.7834} & 0.7668 & 0.7661 \\
Average & 0.7913 & \textbf{0.8076} & 0.8029 & 0.8025 \\
\bottomrule
\end{tabular}
\end{table}

\begin{table}[H]
\centering
\caption{MRPC validation \textit{final} accuracy with a shared epoch-7 budget.}
\label{tab:single-mrpc-fixed}
\small
\begin{tabular}{lcccc}
\toprule
\multirow{2}{*}{Backbone} & baseline r & baseline R & ours r & ours R \\
\midrule
DeBERTa & 0.8520 & \textbf{0.8980} & 0.8907 & 0.8907 \\
RoBERTa & \textbf{0.8912} & 0.8868 & 0.8887 & 0.8877 \\
Average & 0.8716 & \textbf{0.8924} & 0.8897 & 0.8892 \\
\bottomrule
\end{tabular}
\end{table}

\begin{table}[H]
\centering
\caption{SST-2 validation \textit{final} accuracy with a shared epoch-2 budget.}
\label{tab:single-sst2-fixed}
\small
\begin{tabular}{lcccc}
\toprule
\multirow{2}{*}{Backbone} & baseline r & baseline R & ours r & ours R \\
\midrule
DeBERTa & 0.9275 & \textbf{0.9417} & 0.9381 & 0.9372 \\
RoBERTa & 0.9252 & 0.9076 & \textbf{0.9300} & 0.9298 \\
Average & 0.9264 & 0.9247 & \textbf{0.9341} & 0.9335 \\
\bottomrule
\end{tabular}
\end{table}

\begin{table}[H]
\centering
\caption{Task-level averages under the fixed shared-epoch budgets, using
\textit{final} accuracy.}
\label{tab:single-overall-fixed}
\begin{tabular}{lccc}
\toprule
Task & Baseline Avg Final & Ours Avg Final & Margin \\
\midrule
RTE & 0.7995 & \textbf{0.8027} & \textbf{0.0032} \\
MRPC & 0.8820 & \textbf{0.8895} & \textbf{0.0075} \\
SST-2 & 0.9255 & \textbf{0.9338} & \textbf{0.0083} \\
Overall & 0.8690 & \textbf{0.8753} & \textbf{0.0063} \\
\bottomrule
\end{tabular}
\end{table}

Tables~\ref{tab:single-rte-fixed}--\ref{tab:single-overall-fixed} show that
the fixed shared-epoch setting remains favorable to the proposed method on the
aggregate under \textit{final} accuracy. The largest gain appears on MRPC,
SST-2 is also clearly positive, and RTE remains slightly positive.

\begin{table}[H]
\centering
\caption{Worst-case \textit{final} accuracy in the single-task study, defined
as $\min(r, R)$ inside each method family and then averaged over the two
backbones for each task.}
\label{tab:single-worst-final}
\begin{tabular}{lccc}
\toprule
Task & Baseline Worst & Ours Worst & Margin \\
\midrule
RTE & 0.7881 & \textbf{0.8025} & \textbf{0.0144} \\
MRPC & 0.8694 & \textbf{0.8892} & \textbf{0.0199} \\
SST-2 & 0.9175 & \textbf{0.9335} & \textbf{0.0159} \\
Overall & 0.8583 & \textbf{0.8751} & \textbf{0.0167} \\
\bottomrule
\end{tabular}
\end{table}

\begin{table}[H]
\centering
\caption{Overall mean seed variance over the selected single-task windows.
At each kept evaluation point inside the shared-epoch window, we compute
variance across random seeds for the current \textit{final}, the running
\textit{maximum}, and the interpolated
\textit{total}=0.5\cdot\textit{final}+0.5\cdot\textit{maximum}, then average
those variances over the window. Lower is better.}
\label{tab:single-seed-variance}
\begin{tabular}{lccc}
\toprule
Method & Final Var & Max Var & Total Var \\
\midrule
ours r     & \textbf{0.000355} & \textbf{0.000344} & \textbf{0.000340} \\
ours R     & 0.000364 & 0.000351 & 0.000348 \\
baseline r & 0.001549 & 0.001421 & 0.001472 \\
baseline R & 0.000677 & 0.000662 & 0.000662 \\
\bottomrule
\end{tabular}
\end{table}

\begin{table}[H]
\centering
\caption{Supplementary maximum-score view.}
\label{tab:supp-maximum}
\resizebox{\textwidth}{!}{%
\begin{tabular}{lccccccc}
\toprule
Setting & baseline avg maximum & CAGrad avg maximum & ours avg maximum & Margin vs baseline & Margin vs CAGrad & Notes \\
\midrule
Multi mix3 & 0.8174 & 0.8186 & \textbf{0.8230} & \textbf{0.0056} & 0.0044 & ours best \\
Multi mix4 & \textbf{0.8421} & 0.8356 & 0.8408 & -0.0014 & \textbf{0.0052} & baseline best \\
Multi Overall & 0.8298 & 0.8271 & \textbf{0.8319} & \textbf{0.0021} & 0.0048 & ours best \\
Single RTE & \textbf{0.8114} & - & 0.8081 & -0.0032 & - & baseline best \\
Single MRPC & 0.8863 & - & \textbf{0.8937} & \textbf{0.0075} & - & ours best \\
Single SST-2 & 0.9269 & - & \textbf{0.9363} & \textbf{0.0093} & - & ours best \\
Single Overall & 0.8749 & - & \textbf{0.8794} & \textbf{0.0045} & - & ours best \\
\bottomrule
\end{tabular}%
}
\end{table}

\begin{table}[H]
\centering
\caption{Supplementary total-score view. For each retained budget, we define
\textit{total}=0.5\cdot\textit{final}+0.5\cdot\textit{maximum}.}
\label{tab:supp-total}
\resizebox{\textwidth}{!}{%
\begin{tabular}{lccccccc}
\toprule
Setting & baseline avg total & CAGrad avg total & ours avg total & Margin vs baseline & Margin vs CAGrad & Notes \\
\midrule
Multi mix3 & 0.8120 & 0.8154 & \textbf{0.8204} & \textbf{0.0084} & 0.0050 & ours best \\
Multi mix4 & 0.8403 & 0.8337 & \textbf{0.8405} & \textbf{0.0002} & 0.0068 & ours best \\
Multi Overall & 0.8262 & 0.8246 & \textbf{0.8305} & \textbf{0.0043} & 0.0059 & ours best \\
Single RTE & 0.8054 & - & \textbf{0.8054} & \textbf{0.0000} & - & tie on average \\
Single MRPC & 0.8841 & - & \textbf{0.8916} & \textbf{0.0075} & - & ours best \\
Single SST-2 & 0.9262 & - & \textbf{0.9350} & \textbf{0.0088} & - & ours best \\
Single Overall & 0.8719 & - & \textbf{0.8774} & \textbf{0.0054} & - & ours best \\
\bottomrule
\end{tabular}%
}
\end{table}

For the single-task variance metric in Table~\ref{tab:single-seed-variance}, we
do not report only the endpoint variance. Instead, we average seed variance
over the full shared-budget window up to the selected cutoff, using the
current \textit{final}, the running \textit{maximum}, and the interpolated
\textit{total} at each evaluation point.

\section{Appendix: Single-Pass Streaming Accumulation}
\label{app:kernel}

This appendix describes a practical streaming kernel for accumulating all diagnostic consistency statistics
$(C_{r,b}, C_{R,b}, A_b)$ without materializing any vote vectors in memory.
The design explicitly respects the distinct parameter layouts of the $(r)$ and $(hi)$ tensors
under the nested dual-rank LoRA parameterization.

\paragraph{Scope.}
This kernel computes \emph{only} the diagnostic statistics required by the routing rules:
internal consistencies $C_r, C_R$ and the cross-resolution retention ratio $A$.
The statistic $A_b$ is computed directly from the current step's votes.
It does \emph{not} include:
(i) any gating or routing logic,
or (ii) any execution-layer direction construction or gradient aggregation.
In particular, this kernel never produces update directions and does not interact with the optimizer.

We include this kernel as an implementation note for memory-bounded training;
the proposed method itself is independent of this accumulation strategy.

\begin{algorithm}[H]
\caption{Streaming Accumulation Kernel (Per-Block, Unified Voter Set $\mathcal{V}$)}
\begin{algorithmic}[1]
\Require Vote-grad shards on LoRA parameters for block $b$:
$\{g_{U_r}^{(k)}, g_{V_r}^{(k)}, g_{U_{hi}}^{(k)}, g_{V_{hi}}^{(k)}\}_{k\in\mathcal{V}}$.
\Ensure Accumulators $\{S_r^{(k)},D_r^{(i,j)}\}$, $\{S_R^{(k)},D_R^{(i,j)}\}$, and alignment scalars $(X_{\mathrm{sum2}}, Y_{\mathrm{sum2}})$.

\State Initialize all accumulators to $0$.
\State Initialize $X_{\mathrm{sum2}}\gets 0$, $Y_{\mathrm{sum2}}\gets 0$.
\Comment{$X_{\mathrm{sum2}}=\|\sum_{k\in\mathcal V} v_{r,b}^{(k)}\|^2$ over $(U_r,V_r)$; $Y_{\mathrm{sum2}}=\|\sum_{k\in\mathcal V} v_{hi,b}^{(k)}\|^2$ over $(U_{hi},V_{hi})$}

\Statex
\State \textbf{1) Stream over rank-$r$ tensors} $(U_r,V_r)$:
\For{each element index $idx$ in tensors $(U_r,V_r)$}
    \State Load $x_k \gets g_r^{(k)}[idx]$ for $k\in\mathcal{V}$
    \For{$k\in\mathcal{V}$}
        \State $S_r^{(k)} \gets S_r^{(k)} + x_k^2$
        \State $S_R^{(k)} \gets S_R^{(k)} + x_k^2$ \Comment{$r$-part contributes to $R$-view}
    \EndFor
    \For{$i<j,\; i,j\in\mathcal{V}$}
        \State $D_r^{(i,j)} \gets D_r^{(i,j)} + x_i x_j$
        \State $D_R^{(i,j)} \gets D_R^{(i,j)} + x_i x_j$
    \EndFor
    \State $s \gets \sum_{k\in\mathcal V} x_k$
    \State $X_{\mathrm{sum2}} \gets X_{\mathrm{sum2}} + s^2$
\EndFor

\Statex
\State \textbf{2) Stream over rank-extension tensors} $(U_{hi},V_{hi})$:
\For{each element index $idx$ in tensors $(U_{hi},V_{hi})$}
    \State Load $y_k \gets g_{hi}^{(k)}[idx]$ for $k\in\mathcal{V}$
    \For{$k\in\mathcal{V}$}
        \State $S_R^{(k)} \gets S_R^{(k)} + y_k^2$ \Comment{$hi$-part contributes only to $R$-view}
    \EndFor
    \For{$i<j,\; i,j\in\mathcal{V}$}
        \State $D_R^{(i,j)} \gets D_R^{(i,j)} + y_i y_j$
    \EndFor
    \State $t \gets \sum_{k\in\mathcal V} y_k$
    \State $Y_{\mathrm{sum2}} \gets Y_{\mathrm{sum2}} + t^2$
\EndFor

\Statex
\State \textbf{3) Compute statistics:}
\State Compute $C_r$ and $C_R$ from $\{S_r,D_r\}$ and $\{S_R,D_R\}$:
\[
C_{\chi,b}=
\frac{2}{|\mathcal V|(|\mathcal V|-1)}
\sum_{i<j,\; i,j\in\mathcal V}
\frac{D_{\chi}^{(i,j)}}{\sqrt{S_{\chi}^{(i)}}\sqrt{S_{\chi}^{(j)}}+\epsilon},
\quad \chi\in\{r,R\}.
\]
\State Compute the retention ratio using the padding identity:
\[
A_b=
\frac{X_{\mathrm{sum2}}}
{\sqrt{X_{\mathrm{sum2}}}\sqrt{X_{\mathrm{sum2}}+Y_{\mathrm{sum2}}}+\epsilon}.
\]
\end{algorithmic}
\end{algorithm}

\paragraph{Scalar storage (per block, variable $V_t$).}
\begin{center}
\begin{tabular}{lccc}
\toprule
Quantity & Accumulators & Count & Storage (floats) \\
\midrule
$r$-Consistency ($C_r$) & $S_r, D_r$ & $V_t + \tfrac{V_t(V_t-1)}{2}$ & $V_t + \tfrac{V_t(V_t-1)}{2}$ \\
$R$-Consistency ($C_R$) & $S_R, D_R$ & $V_t + \tfrac{V_t(V_t-1)}{2}$ & $V_t + \tfrac{V_t(V_t-1)}{2}$ \\
Retention Ratio ($A_b$) & $X_{\mathrm{sum2}}, Y_{\mathrm{sum2}}$ & $2$ & $2$ \\
\midrule
\textbf{Total (diagnostic scalars)} & & & $\mathbf{V_t^2 + V_t + 2}$ \\
\bottomrule
\end{tabular}
\end{center}

\paragraph{Compute cost.}
Per-parameter cost is $O(|\mathcal V|^2)$ due to pairwise dot accumulators.
This trades additional compute for bounded memory by avoiding
materialization of vote vectors.

\bibliographystyle{plain}
\bibliography{references}
\end{document}
